\documentclass[11pt,a4paper]{article}
\usepackage[margin=0.85in,top=1in,bottom=1in]{geometry}
\usepackage{amsmath,amssymb,amsfonts}
\usepackage{graphicx}
\usepackage{float}
\usepackage{enumitem}
\usepackage{microtype}
\usepackage{caption}
\usepackage[hidelinks]{hyperref}
\usepackage[most,breakable]{tcolorbox}
\tcbuselibrary{breakable,skins}
\usepackage[utf8]{inputenc}
\usepackage[T1]{fontenc}
\usepackage{lmodern}
\captionsetup{font=small,labelfont=bf,skip=4pt}
\setlist{leftmargin=1.5em,topsep=2pt}

%% Breakable card — prevents content cutoff across pages
\newtcolorbox{solcard}[3]{
  breakable, enhanced,
  colback=white, colframe=gray!50, arc=2pt,
  title={\textbf{#1}\hfill\textsf{\small #3\,pts}},
  fonttitle=\normalsize\sffamily\bfseries,
  before upper={\small\textit{#2}\par\smallskip\hrule height 0.4pt\smallskip},
  top=4pt, bottom=6pt, left=7pt, right=7pt,
  parbox=false,
}

\title{\textbf{Electrons in a magnetic field --- Solution}}
\author{\normalsize Y20 (2020) — English translation}
\date{}
\begin{document}\maketitle\vspace{-0.5em}
\begin{solcard}{A1}{What could be the minimum distance $d_{\min}$ between electrons so that the movement described above is possible?}{1.00}

Let's write Newton's second law in projection onto a line passing through the electrons. Since the distance does not change, the acceleration is $\frac{v^2}{r}$. We end up with the equation: $$ \frac{2mv^2}{d}=vBe-\frac{e^2}{4\pi\varepsilon_0 d^2} $$


We solve the quadratic equation for $v$. The discriminant of this equation is $(Be)^2-\frac{2me^2}{\pi\varepsilon_0d^3}$.


From the discriminant it is clear that solutions for speed exist only for $d\geqslant2\sqrt[3]{\frac{m}{4\pi\varepsilon_0 B^2}}$. This will be the minimum possible distance.

\end{solcard}

\begin{solcard}{A2}{Find the speed $v_m$ that must be imparted to the electrons so that when they move between them, the distance $d_{\min}$ found in A1 is maintained.}{0.50}

Since the discriminant is zero: $$v_m = -\frac{b}{2a}= \frac{Be}{2m}d_{\text{min}}=\frac{e}{2}\sqrt[3]{\frac{B}{4\pi\varepsilon_0m^2}}$$

\end{solcard}

\begin{solcard}{B1}{Find the angular velocity of the center of mass $\omega$ during the electron motion.}{0.50}

The center of mass moves under the influence of two Lorentz forces: $$m\vec{a} = \left[\vec{v}_1\times \vec B\right]q+\left[\vec{v}_2\times \vec B\right]q = \left[\left(\vec{v}_1+\vec{v}_2\right)\times \vec B\right]q = 2\left[\vec{v}_{CM}\times \vec B\right]q$$


Since there are no other external forces, the speed of the center of mass rotates, but does not change in magnitude. The acceleration of the center of mass in this case is equal to $\vec\omega\times\vec v$ $$2m\left[\vec\omega_{CM}\times\vec v\right] \equiv 2\left[\vec{v}_{CM}\times \vec B\right]q$$ $$\omega_{CM} = \frac{Be}{m}$$

\end{solcard}

\begin{solcard}{B2}{Find the time it takes for both electrons to have the same absolute speed for the first time.}{1.00}

During the movement described in the condition, kinetic energy is conserved. Thus, the angular velocity relative to the center of mass does not change. $$\omega = \frac{v_0}{d_1}$$


The speed of the center of mass and the relative speed are equal, so the speeds of the electrons will become equal when the angle between the speed of the center of mass and the relative speed becomes equal to $\frac{\pi}{2}$. $$t = \frac{\pi}{2}\frac{1}{\left|\frac{v_0}{d_1} - \frac{eB}{m}\right|}$$

\end{solcard}

\begin{solcard}{C1}{Obtain an expression for the kinetic energy of one electron $K=K(m, v_r, \omega, r)$, where $v_r=\cfrac{dr}{dt}, \omega=\cfrac{d\theta}{dt}$, and $r$ and $\theta$ specify the position of the electron in the center of mass system.}{0.20}

$K=\frac{m}{2}(\omega^2r^2+v_r^2)$

\end{solcard}

\begin{solcard}{С2}{Find what the time derivative of angular momentum $\cfrac{dL}{dt}$ is equal to when electrons move.}{0.80}

The angular momentum changes due to the Lorentz force component corresponding to the radial velocity component. $$\frac{dL}{dt} = v_rBer$$

\end{solcard}

\begin{solcard}{C3}{Find the integral of motion for the case under consideration, i.e. find a function $J$ for which $\cfrac{dJ}{dt}=0$ is true.}{0.50}

Integrating the expression for $\frac{dL}{dt}$ we obtain the integral of motion $$J = L-\frac{eBr^2}{2}$$

\end{solcard}

\begin{solcard}{C4}{Express the total energy of the electron as $E=K(v_r)+U(r)$. Express it in terms of $J, e, B, r$ and physical constants.}{1.00}

We express the angular velocity through $J$, then substitute it into the expression for kinetic energy. $\omega  = \frac{1}{m}\left(\frac{J}{r^2}+\frac{eB}{2}\right)$.


To the kinetic energy we must add the energy of the Coulomb interaction per electron. $$E = \frac{mv_r^2}{2}+\frac{r^2}{2m}\left(\frac{J}{r^2}+\frac{eB}{2}\right)^2+\frac{e^2}{16\pi\varepsilon_0r}$$

\end{solcard}

\begin{solcard}{C5}{Schematically depict the graph of the dependence $U(r)$, indicate its characteristic features.}{0.50}

The limits of $U(r)$ at points $0$ and $+\infty$ are equal to $+\infty$, and the function is continuous on the segment $(0, +\infty)$


The second derivative of the function $U(r)$ has the following form: $$\frac{Ar^4+Br+C}{r^4}; ~~A,~B,~C, ~r > 0$$ On the segment of interest to us, it is always greater than zero, both limits are infinite, therefore there is a minimum, and the only one.

\begin{figure}[H]
  \centering
  \includegraphics[width=0.72\textwidth]{T5_S_EN_assets/fig1.jpg}
\end{figure}
\end{solcard}

\begin{solcard}{С6}{Expand the function $U(r)$ in a Taylor series up to the first non-trivial term near the singular points from C5. Based on the decomposition, sketch the trajectory of electrons that have been given a small radial velocity component $v_{r0}$ in addition to the circular velocity $v_0$.}{1.00}

Since we will remain in the vicinity of the minimum, the electrons will perform harmonic oscillations along the radius.

\begin{figure}[H]
  \centering
  \includegraphics[width=0.72\textwidth]{T5_S_EN_assets/fig2.jpg}
\end{figure}
\end{solcard}

\end{document}